% Part: first-order-logic
% Chapter: syntax-and-semantics
% Section: introduction

\documentclass[../../../include/open-logic-section]{subfiles}

\begin{document}

\olfileid[pt]{fol}{syn}{int}

\olsection{Introdução}

Para desenvolver a teoria e a metateoria da lógica de primeira ordem,
devemos primeiro definir a sintaxe e a semântica de suas expressões.  As
expressões da lógica de primeira ordem são termos e !!{formula}s.  Os
termos são formados a partir de !!{variable}s, !!{constant}s e
!!{function}s. As !!^{formula}s, por sua vez, são formadas a partir de
!!{predicate}s juntamente com termos (estes formam as !!{formula}s
atômicas, as menores), e então a partir de !!{formula}s atômicas
podemos formar outras mais complexas usando conectivos lógicos e
quantificadores.  Há muitas maneiras diferentes de estabelecer as
regras de formação; apresentamos apenas uma possível. Outros sistemas
escolherão símbolos diferentes, selecionarão conjuntos diferentes de
conectivos como primitivos, usarão parênteses de maneira diferente (ou
nem os usarão, como no caso da chamada notação polonesa).  O que todas
as abordagens têm em comum, no entanto, é que as regras de formação
definem o conjunto de termos e !!{formula}s \emph{indutivamente}. Se
feito corretamente, toda expressão pode resultar essencialmente de
apenas uma maneira de acordo com as regras de formação.  A definição
indutiva que resulta em expressões de leitura \emph{única} significa
que podemos dar significados a essas expressões usando o mesmo
método---definição indutiva.

Dar o significado das expressões é o domínio da semântica.  O conceito
central na semântica é o de satisfação em !!a{structure}. Uma
!!^a{structure} dá significado aos blocos de construção da linguagem:
um !!a{domain} é um conjunto não vazio de objetos. Os quantificadores
são interpretados como variando sobre esse domínio, !!{constant}s são
atribuídas a elementos do domínio, !!{function}s são atribuídas a
funções do !!{domain} em si mesmo, e !!{predicate}s são atribuídos a
relações sobre o !!{domain}.  O !!{domain} juntamente com as
atribuições ao vocabulário básico constitui !!a{structure}. As
!!^{variable}s podem aparecer em !!{formula}s e, para dar uma semântica,
também temos que atribuir !!{element}s do !!{domain} a elas---isto é
uma atribuição de variáveis. A relação de satisfação, finalmente, une
tudo isso. Uma !!^a{formula} pode ser satisfeita em !!a{structure}
~$\Struct{M}$ em relação a uma atribuição de variáveis~$s$, escrito como
$\Sat{M}{!A}[s]$. Esta relação também é definida por indução na
estrutura de~$!A$, usando as tabelas-verdade dos conectivos lógicos
para definir, por exemplo, a satisfação de $!A \land !B$ em termos da
satisfação (ou não) de $!A$ e de~$!B$.  Verifica-se então que a
atribuição de variáveis é irrelevante se a !!{formula}~$!A$ é
!!a{sentence}, isto é, não tem variáveis livres, de modo que podemos
falar de !!{sentence}s sendo simplesmente satisfeitas (ou não) em
!!{structure}s.

Com base na relação de satisfação $\Sat{M}{!A}$ para sentenças,
podemos então definir as noções semânticas básicas de validade,
consequência e satisfatibilidade.  Uma sentença é válida, $\Entails !A$,
se toda estrutura a satisfaz. Ela é consequência de um conjunto de
!!{sentence}s, $\Gamma \Entails !A$, se toda !!{structure} que
satisfaz todas as !!{sentence}s em~$\Gamma$ também satisfaz~$!A$. E um
conjunto de sentenças é satisfatível se alguma !!{structure} satisfaz
todas as !!{sentence}s nele ao mesmo tempo.  Como as !!{formula}s são
definidas indutivamente, e a satisfação é, por sua vez, definida por
indução na estrutura das !!{formula}s, podemos usar indução para provar
propriedades da nossa semântica e relacionar as noções semânticas
definidas.




\end{document}
